\section{An autonomous discovery workflow and the limits of its authority}
\label{exec17:copilot-development}

An analyst who can run a discovery algorithm still faces a substantial sequence
of choices: expressing the question, checking the observation clock, choosing
an admissible method, configuring it, interpreting its graph and communicating
what the graph permits. Causal-Copilot addresses much of this sequence in a
single agent workflow. Wang and colleagues describe natural-language input,
data diagnostics, algorithm selection and configuration, execution revision,
graph refinement and report generation. This is a useful architectural
comparison for the proposed banking assistant. Its value lies partly in
organizing work that a numerical discovery routine leaves to the analyst.
The extent of that automation and the causal authority of its output require
separate examination.\citep{exec17_copilot_paper}

The evidence considered here has two distinct versions. The paper is
arXiv:2504.13263v2, dated 21 April 2025. The inspected repository snapshot is
commit \texttt{a31aa7c}, dated 15 May 2026. Both precede the literature cutoff,
but the later source has not been established as the code that produced the
paper's experiments. Paper results below are author-reported results;
source findings describe particular static paths in the later snapshot.
Neither supplies a reproduction of the other. This distinction permits a
useful component assessment without attributing a later implementation choice
to an earlier experiment.\citep{exec17_copilot_paper,exec17_copilot_code}

\subsection{Translating a retention question into a causal request}

Suppose a manager asks, ``Which retention offers prevent withdrawals?''
Causal-Copilot's input and preprocessing modules offer a concrete starting
point: parse a natural-language objective and constraints, inspect the supplied
data, and allow the user to revise the request. Its subsequent modules filter
and rank algorithms, configure execution and revise errors. The paper also
describes an interactive graph revision loop and a generated PDF report;
its web interface is an author-described interface rather than a deployment
tested here. These capabilities show why the proposal needs a complete
analyst journey, including intelligible reasons for method selection, rather
than a collection of numerical commands.\citep{exec17_copilot_paper}

For the retention request, a proposed causal contract must go further than
extracting the words ``offer'' and ``withdrawal.'' Let $X$ denote assignment of
a specified offer before a decision deadline, and let $Y$ indicate retention
over a fixed subsequent horizon. The population must specify eligibility at
that deadline. A receipt date, an acceptance date and an assignment date can
produce different treatments; an account balance observed after an offer can
be a consequence of treatment rather than an admissible adjustment variable.
Multiple offers require either a specified version or an explicit policy.
The assistant should expose these alternatives and ask for the information
that determines the intended contrast. A fluent question-to-graph translation
does not establish the timing or truth of the underlying variables.

The selected source makes part of the selection process concrete. A user-chosen
algorithm bypasses the filtering step. Otherwise, the filter sends the query,
data statistics, supplied knowledge, resource information and the requested
DAG or CPDAG form to a temperature-zero language-model selection step. Such a
prompt can organize alternatives and enforce an output preference. It cannot
make an observationally undetermined direction identifiable merely because a
DAG was requested. The proposal should retain the candidate class, rejected
alternatives and assumptions that justify the selection. User authority to
choose an algorithm is likewise distinct from evidence that its assumptions
hold.\citep{exec17_copilot_code}

The proposed contract would therefore distinguish a documented business fact,
a modeling assumption, a statistical diagnostic and an agent suggestion.
For example, an audited assignment rule may justify an independence premise;
an apparent lack of linear association cannot certify it. A suggested latent
``relationship strength'' variable may explain selection, but naming it does
not measure it. In the eventual CLI or MCP request, these distinctions need
structured fields that survive the call to a numerical engine and its returned
report. They are proposed additions to this library, with semantic and runtime
qualification still required.

\subsection{Stable association can support incompatible decisions}

The distinction between workflow assistance and identification has a small
exact demonstration. Let $U$ and $V$ be independent, with
$U\sim\operatorname{Bernoulli}(1/2)$ and
$V\sim\operatorname{Bernoulli}(1/10)$. Write $\oplus$ for binary exclusive-or.
Consider two stipulated structural worlds:
\begin{equation}
 \begin{array}{lll}
 A:& X=U,&Y=X\oplus V,\\
 B:& X=U\oplus V,&Y=U.
 \end{array}
 \label{exec17:copilot-worlds}
\end{equation}
In world $A$, the offer changes retention. In world $B$, the hidden customer
state $U$ precedes both assignment and the later outcome, producing selection
without an offer effect. These are mathematical examples, not observed banking
mechanisms or fitted Causal-Copilot models. Both induce the same strictly
positive observed distribution:
\begin{equation}
 P(0,0)=P(1,1)=\frac9{20},\qquad
 P(0,1)=P(1,0)=\frac1{20}.
 \label{exec17:copilot-observed}
\end{equation}
Consequently, $E(Y\mid X=1)-E(Y\mid X=0)=4/5$ in either world.
An intervention replacing the equation for $X$ leaves $Y=U$ unchanged in
world $B$. In world $A$, it gives retention probabilities $9/10$ and $1/10$.
The interventional contrast is therefore $4/5$ in $A$ and zero in $B$.
Enumerating the four $(U,V)$ atoms reproduces both the observed law and these
effects. The pair demonstrates nonidentification in a class admitting both
worlds; it is neither a general sharp effect bound nor a statement about an
actual discovery run.\citep{exec17_copilot_exact}

Now condition on six independent observations comprising three $(0,0)$ pairs
and three $(1,1)$ pairs. In an ordinary iid empirical bootstrap of size six,
the covariance of $X$ and $Y$ is positive whenever both types appear.
It vanishes only for either of the two constant-type resamples. Hence
\begin{equation}
 P_*\{\widehat{\operatorname{Cov}}(X,Y)>0\}
 =1-2(1/2)^6=\frac{31}{32}=0.96875.
 \label{exec17:copilot-stability}
\end{equation}
Here $P_*$ is conditional resampling probability, not probability that a causal
direction is true. Exhausting all $6^6=46{,}656$ index resamples gives 45,198
positive covariances. The conditioning event has probability
$\binom63(9/20)^6=531441/3200000$ under either observed law, so it is a possible
sample with an explicit probability. The stability calculation concerns a
covariance sign; it does not claim that any specific upstream algorithm returns
an edge with that frequency.\citep{exec17_copilot_exact}

With a stipulated USD\,100 contribution per incremental retained customer and
a USD\,30 expense per assigned offer, the contrast against withholding is
\begin{equation}
 \Delta_A=100(4/5)-30=50,\qquad
 \Delta_B=100(0)-30=-30\quad\text{USD per offer}.
 \label{exec17:copilot-value}
\end{equation}
The same observed data can thus favor opposite actions in the two admitted
worlds. More confident repetition of the association does not select between
them. A useful assistant should identify what would change this conclusion:
a credible assignment experiment, defensible additional restrictions, or new
information that excludes a world. Which information is feasible remains a
banking design question. Bootstrap stability is still useful for describing
sampling sensitivity under a declared resampling law; its interpretation must
stop at that law.\citep{exec17_copilot_exact}

\subsection{What the reported discovery comparisons establish}

The paper's synthetic experiments supply known graphs and vary the generating
conditions. They provide evidence about selection and discovery within those
conditions, including favorable settings and substantive alternatives. The
default tabular row in Table 2 reports Causal-Copilot F1 of 0.900, compared with
0.920 for PC, 0.910 for FCI and 0.920 for GES. In the sparse row, its 0.820 is
below GPT-4o's 0.830 and FCI's 0.840. These countercomparisons matter when
evaluating an adaptive selector: automation can be useful without dominating
each fixed method. They do not support a universal superiority claim.
\citep{exec17_copilot_paper}

The financial scenario is generated, with 50 nodes across three financial
domains; the reported Copilot F1 is $0.650\pm0.130$. Financial vocabulary does
not turn this known-truth simulation into an authentic bank validation. The
temporal Table 3 is explicitly linear and stationary. Its default row gives
0.673 for Copilot, 0.695 for PCMCI and 0.733 for DYNOTEARS. The table's
$n=5000$ row reports 0.668 for Copilot, whereas the adjacent prose says 0.649.
The very-large row labels $p=100$, whereas part of the prose says $p=1000$.
These two disagreements remain unresolved; no raw output or experiment-bound
configuration is supplied by this assessment to reconcile them.
\citep{exec17_copilot_paper}

The evaluation description on physical page 47 imposes a twenty-minute runtime
cap, describes a best-DAG conversion for tabular predictions and evaluates
temporal summary graphs while excluding diagonal entries. Such conventions
change the evaluated object. A summary adjacency can record that some lagged
relationship exists while losing which lag matters. Excluding a diagonal
also excludes a variable's own autoregressive relationships from that graph
score. Completion under a runtime cap is relevant operational evidence, but
an unavailable run and an incorrect completed graph are different outcomes.
A banking comparison should retain both, specify the scored graph object,
and examine the intended intervention query rather than infer effect quality
from a summary F1 alone.\citep{exec17_copilot_paper}

The later source supplies a precise example of why the conversion convention
needs an explicit name. In \texttt{evaluator.py}, entries $[i,j]=1$ mean
$j\to i$, and the default is \texttt{sample=False}. The default metric method
calls a helper that consults \texttt{true\_graph} when orienting ambiguous
predicted edges. Its alternative sampling branch also selects using minimum
SHD against truth. The integrated tabular benchmark constructs the default
evaluator and reaches this helper; its temporal branch instead calls direct
binary scoring. This is truth-assisted benchmark evaluation in the inspected
paths. It is not evidence that a deployed discovery procedure receives the
unknown true graph, and it has not been bound to the 2025 experiments.
\citep{exec17_copilot_code}

For two nodes, hold fixed a single legal asymmetric undirected prediction and
consider opposite one-edge true DAGs. Both directions are compatible with this
unoriented prediction. With ordered off-diagonal binary F1, a truth-assisted
conversion can attain one in both worlds. A conversion fixed to one direction
before seeing truth attains one in one world and zero in the other. Under a
stipulated equal weighting of the worlds,
\begin{equation}
 \overline{F1}_{\mathrm{assisted}}=1,\qquad
 \overline{F1}_{\mathrm{fixed}}=\frac12.
 \label{exec17:copilot-score}
\end{equation}
This authored enumeration isolates the information used by the score; it is
not an estimate of the paper's score inflation. An optimistic compatibility
score can answer a legitimate question about whether an ambiguous output
contains the truth. It should be reported as that question, alongside ambiguity
and a declared prediction rule, rather than interpreted as a direction the
algorithm selected without truth.\citep{exec17_copilot_exact}

\subsection{A graph-refinement invariant needs an executable guard}

Physical page 7 of the paper states that LLM recommendations are soft
suggestions and may not override high-confidence bootstrap decisions. That
is an intelligible division of authority. A source-level assessment must ask
whether the reachable editing path enforces it, including the representation
of opposite directions of the same pair. In the later snapshot,
\texttt{bootstrap\_recommend} identifies high directed frequencies at 0.9 or
above and low frequencies at 0.1 or below. Low ordered entries also enter its
``middle'' collection. The LLM evaluation combines that collection and
deduplicates it as unordered pairs. Its prompt distinguishes two directions,
answer C for no relationship and answer D for unknown.
\citep{exec17_copilot_paper,exec17_copilot_code}

Consider a conditional static witness using the source's receiver/source
matrix convention. Stipulate a raw $0\to1$ edge and twenty bootstrap results,
nineteen in that direction and one in the reverse direction:
\begin{equation}
 G=\begin{pmatrix}0&0\\1&0\end{pmatrix},\qquad
 B=\begin{pmatrix}0&0.05\\0.95&0\end{pmatrix}.
 \label{exec17:copilot-edit-input}
\end{equation}
Set uncertain-edge frequencies to zero. The forward ordered entry is high;
the reverse entry is low and therefore also enters the middle collection.
Unordered deduplication makes the pair eligible for LLM examination. In the
inspected \texttt{quality\_judge} path, \texttt{fixed\_pairs} is empty. If an
answer C is accepted for this pair, the resulting forbid entry reaches an edit
that sets both entries of $G$ to zero. A separately authored finite translation
of these expressions reproduces the mutation.
\citep{exec17_copilot_code,exec17_copilot_exact}

The conditions in this witness are essential. The bootstrap collection and
accepted answer are stipulated, not observed model outputs. The main program
reaches the judge and its quality path, but no upstream process or background
knowledge implementation was executed here. The witness stops immediately
after the LLM mutation; later KCI pruning and cycle handling are not simulated.
It therefore identifies an unguarded intermediate edit under stated inputs,
rather than claiming a reproduced final graph failure. It also does not
establish that the 2025 implementation violated the paper's rule.
\citep{exec17_copilot_code,exec17_copilot_exact}

A proposed repair would encode protection on unordered pairs before any edit:
\begin{equation}
 \begin{gathered}
 \mathcal P=\{\{i,j\}:i\ne j,\ B_{ij}\ge0.9\},\\
 \{i,j\}\in\mathcal P\ \Longrightarrow\
 G^{\mathrm{after\ LLM}}_{\{i,j\}}=G^{\mathrm{before\ LLM}}_{\{i,j\}}.
 \end{gathered}
 \label{exec17:copilot-guard}
\end{equation}
The equality refers to both ordered entries for the protected pair. Applying
this rule blocks the same conditional deletion in the authored reference.
A full integration would have to define the precedence of user changes,
statistical constraints and subsequent pruning, then verify the invariant at
each relevant boundary. Protecting an edge implements an authority rule; it
does not prove that the protected edge is causally correct. The two-world
example already shows why such a proof cannot follow from frequency alone.
\citep{exec17_copilot_exact}

\subsection{Temporal semantics, runtime scope and component adoption}

Temporal support requires following the path that is actually called. The
inspected bootstrap helper has an iid row branch and a temporal branch with
block size ten and block starts sampled without replacement. That is a
particular resampling scheme, not automatically a justified time-series
bootstrap. More directly, the inspected judge calls the helper with
\texttt{ts=False}, even though the helper offers a temporal branch. The
program wrapper can preserve returned lag matrices, and the state declares
raw, revised and lagged fields. Storage capacity does not establish that the
judge resamples, edits and reports the same temporal object.
\citep{exec17_copilot_code}

For a servicing query, an edge $X_{t-1}\to Y_t$ and an edge
$X_{t-5}\to Y_t$ imply different feasible intervention times. Summarizing
both as $X\to Y$ can be useful for an overview while being insufficient for
the requested policy. A proposal integration should preserve lag, endpoint
type, observation clock and the declared summary operation in its result.
It should select a resampling law justified for that target and explicitly
refuse a temporal confidence claim when only an iid route is available.
Whether a fixed block size and without-replacement scheme have suitable
behavior depends on the assumed dependence and asymptotics; the inspected
branch alone supplies neither a theorem nor a qualification study.

Runtime claims also need a boundary. The inspected benchmark timer encloses
\texttt{Programming.forward}, rather than the whole natural-language,
diagnostic, selection, refinement and reporting journey. Its source also
replaces missing values with zero, a preprocessing choice whose causal
consequences depend on the missingness mechanism. Default pytest settings
exclude integration, LLM and LaTeX markers; this is evidence of declared tests
and selected exclusions, not absence of tests. The requirements contain pinned
entries alongside an unpinned \texttt{huggingface-hub}. An environment lock,
end-to-end runtime measurement and semantic acceptance suite therefore remain
necessary before a reproducible runtime comparison.
\citep{exec17_copilot_code}

The bounded adoption decision is to use the workflow decomposition as an
architectural reference and assess components behind explicit contracts.
Question parsing, method explanations and report drafting are candidates for
adaptation with evidence-linked inputs and human revision. A statistical
wrapper is a candidate for reuse after validating its graph and temporal
semantics against an independent reference. The truth-assisted evaluator is
unsuitable as the sole operational accuracy score; it can be retained only as
a separately named compatibility diagnostic. The unguarded refinement path
requires a repair and invariant checks before adopting the paper's authority
promise. None of these static decisions qualifies an installed package.
\citep{exec17_copilot_code,exec17_copilot_exact}

The root license grants MIT permissions with a 2026 contributor notice; the
inspected vendored causal-learn notice grants MIT permissions with a 2022
notice. Both require preservation of the notice and disclaim warranty. These
two originals support a bounded licensing assessment of the inspected files.
They do not settle every transitive dependency, model provider, data asset,
deployment or redistribution right. Direct adoption remains conditional on
that complete inventory and on a tested environment; a port would create an
additional numerical and semantic equivalence obligation rather than removing
the rights question.\citep{exec17_copilot_rights}

The comparison also exposes deficits in this proposal. It has not demonstrated
an integrated Copilot-like business conversation, algorithm selector and
interactive graph revision experience. Nor has it qualified the proposed
ambiguity-preserving engine boundary, temporal guard, evidence-linked report
or complete workflow costs. These are concrete prototype objectives. A first
public-data assessment can combine known-mechanism simulations with paired
requests that change assignment knowledge, lag information or the authority
to edit a graph. It should retain refusals, failures, ambiguity and total
workflow expense in the denominator, and use the exact examples above as
small reference cases. Such an assessment can justify whether an expensive
bank experiment is worth conducting. It cannot substitute public observational
labels for unknown causal truth or declare a banking product validated.
\label{exec17:copilot-development-end}
